% =====================================================================
\section{Cơ sở lý thuyết}
% =====================================================================

% ---------------------------------------------------------------------
\begin{frame}{Kiến trúc NeoBERT}
\vspace{-0.35em}
\centering
\includegraphics[width=\linewidth,height=7.05cm,keepaspectratio]{fig-neobert-arch}
\end{frame}

% ---------------------------------------------------------------------
\begin{frame}{RoPE: mã hóa vị trí bằng phép quay véc-tơ}
\vspace{0.2em}
\centering
\includegraphics[width=0.84\linewidth,height=5.4cm,keepaspectratio]{fig-rope-paper}

\vspace{0.2em}
{\footnotesize\color{vnCite} Nguồn: \citet{Su2024}}
\end{frame}

% ---------------------------------------------------------------------
\begin{frame}{Pre-RMSNorm: gradient ổn định hơn, phép tính rẻ hơn}
\vspace{-0.35em}
\begin{columns}[T]
% ---------------- cột trái: vị trí đặt chuẩn hóa ----------------
\begin{column}{0.40\textwidth}
\centering
\resizebox{\ifdim\width>\linewidth\linewidth\else\width\fi}{!}{%
\begin{tikzpicture}[
  bx/.style={draw=vnNavy,fill=vnLight,rounded corners=2pt,text width=1.75cm,
             align=center,font=\scriptsize,minimum height=0.5cm,inner sep=2pt},
  nm/.style={bx,fill=white},
  pl/.style={draw=vnNavy,circle,inner sep=1.2pt,font=\scriptsize},
  ar/.style={-{Latex[length=1.4mm]},vnNavy,semithick}]

  % ---- Post-LN ----
  \node[font=\bfseries\scriptsize,text=vnNavy] at (0,4.55) {Post-LN (BERT)};
  \node[font=\scriptsize] (a0) at (0,0)     {$x$};
  \node[bx]               (a1) at (0,1.0)   {Lớp con};
  \node[pl]               (a2) at (0,2.0)   {$+$};
  \node[nm]               (a3) at (0,3.05)  {LayerNorm};
  \node[font=\scriptsize] (a4) at (0,3.95)  {$y$};
  \draw[ar] (a0)--(a1); \draw[ar] (a1)--(a2);
  \draw[ar] (a2)--(a3); \draw[ar] (a3)--(a4);
  \draw[ar] (a0.east) -- ++(1.3,0) |- (a2.east);
  \node[font=\tiny,text=vnCite,align=center] at (0,-0.75)
       {chuẩn hóa nằm \emph{trên}\\residual};

  % ---- Pre-RMSNorm ----
  \begin{scope}[xshift=3.9cm]
    \node[font=\bfseries\scriptsize,text=vnNavy] at (0,4.55) {Pre-RMSNorm};
    \node[font=\scriptsize] (b0) at (0,0)     {$x$};
    \node[nm]               (b1) at (0,1.0)   {RMSNorm};
    \node[bx]               (b2) at (0,2.05)  {Lớp con};
    \node[pl]               (b3) at (0,3.05)  {$+$};
    \node[font=\scriptsize] (b4) at (0,3.95)  {$y$};
    \draw[ar] (b0)--(b1); \draw[ar] (b1)--(b2);
    \draw[ar] (b2)--(b3); \draw[ar] (b3)--(b4);
    \draw[ar] (b0.east) -- ++(1.3,0) |- (b3.east);
    \node[font=\tiny,text=vnCite,align=center] at (0,-0.75)
         {residual \emph{thông suốt}\\từ đầu đến cuối};
  \end{scope}
\end{tikzpicture}}
\end{column}

% ---------------- cột phải: hai lý do ----------------
\begin{column}{0.575\textwidth}
\small
\lead{1. Đặt \emph{trước} $\Rightarrow$ gradient không phình ở lớp cuối}
Chuẩn gradient tại lớp cuối $\|\partial\mathcal{L}/\partial W^{L,2}\|_F$
khi khởi tạo ($L$ lớp, $d$ chiều):
\vspace{-0.35em}
\[\textstyle
  \underbrace{\mathcal{O}\big(d\sqrt{\ln d}\big)}_{\text{Post-LN}}
  \qquad\text{so với}\qquad
  \underbrace{\mathcal{O}\big(d\sqrt{\ln d\,/\,L}\big)}_{\text{Pre-LN}}
\]
\vspace{-0.85em}

Post-LN chuẩn hóa lại residual ở \emph{mọi} lớp: gradient
\textbf{không nhỏ đi} khi mạng sâu thêm, buộc phải dùng
\textit{warm-up} tốc độ học. Pre-LN chia thêm $\sqrt{L}$ nên \hl{28 lớp} của
NeoBERT huấn luyện ổn định \cit{Xiong2020}

\vspace{0.15em}
\lead{2. Bỏ trung bình $\Rightarrow$ mỗi lớp rẻ hơn}
\vspace{-0.45em}
\[\textstyle
  \underbrace{\tfrac{a_i-\mu}{\sigma}g_i+b_i}_{\text{LayerNorm}}
  \;\longrightarrow\;
  \underbrace{\tfrac{a_i}{\mathrm{RMS}(a)}g_i}_{\text{RMSNorm}},\;
  \mathrm{RMS}(a)=\sqrt{\tfrac{1}{n}\textstyle\sum_i a_i^2}
\]
\vspace{-0.5em}
\begin{flushright}\cit{Zhang2019}\end{flushright}
\end{column}
\end{columns}
\end{frame}

% ---------------------------------------------------------------------
\begin{frame}{Ba đặc điểm của NeoBERT định hình cách chuyển giao}
\vspace{0.4em}
\begin{columns}[T]
\begin{column}{0.48\textwidth}
  \lead{1. RoPE -- thuận lợi}
  \begin{itemize}\itemsep2pt
    \item Vị trí được áp \emph{bên trong} mỗi lớp, \textbf{không nằm} trong
          ma trận véc-tơ.
    \item[$\Rightarrow$] Thay toàn bộ lớp từ vựng tiếng Việt
          \emph{không đụng chạm} cơ chế vị trí.
  \end{itemize}
\end{column}

\begin{column}{0.48\textwidth}
  \lead{2. Pre-RMSNorm -- ràng buộc}
  \begin{itemize}\itemsep2pt
    \item Chuẩn hóa chia theo \emph{độ lớn} véc-tơ; lớp véc-tơ đi thẳng vào
          luồng chính qua residual.
    \item[$\Rightarrow$] Nếu độ lớn véc-tơ mới lệch khỏi thống kê NeoBERT,
          sai lệch bị khuếch đại qua 28 lớp.
  \end{itemize}
\end{column}
\end{columns}

\vspace{0.9em}
\begin{alertblock}{3. Lớp đầu ra không dùng chung trọng số}
  Khác BERT/RoBERTa, đầu MLM của NeoBERT là một ma trận \emph{độc lập} với
  ma trận véc-tơ đầu vào. Vì vậy chuyển sang tiếng Việt phải khởi tạo
  \hl{cả hai ma trận}, mỗi ma trận trong đúng không gian của nó.
\end{alertblock}
\end{frame}

% ---------------------------------------------------------------------
\begin{frame}{Đổi ngôn ngữ: khởi tạo lớp véc-tơ thế nào?}
\vspace{0.4em}
\centering
% Diagram is narrower than the text block, so \fitpic would leave it small.
% Force it up to 0.86\linewidth instead.
\resizebox{0.93\linewidth}{!}{\begin{tikzpicture}[font=\scriptsize,
  b/.style={draw=vnNavy,rounded corners=2pt,text width=2.6cm,
            minimum height=16mm,align=center,inner sep=3pt},
  ar/.style={-{Latex[length=2mm]},vnSteel,thick}]

  \node[b,fill=vnRed!10,draw=vnRed] (rnd)
       {\textbf{Ngẫu nhiên}\\[-1pt]\tiny tín hiệu vô nghĩa};
  \node[b,fill=vnLight,right=9mm of rnd] (sem)
       {\textbf{Theo ngữ nghĩa}\\[-1pt]\tiny WECHSEL, FOCUS, OFA};
  \node[b,fill=vnNavy,text=white,draw=vnNavy,right=9mm of sem] (salt)
       {\textbf{SALT}\\[-1pt]\tiny tái sử dụng không gian};

  \draw[ar] (rnd)--(sem);
  \draw[ar] (sem)--(salt);

  \node[below=3mm of rnd,text=vnRed,font=\tiny,align=center]
       {phá vỡ biểu diễn\\ đã có};
  \node[below=3mm of sem,text=vnCite,font=\tiny,align=center]
       {mượn \emph{véc-tơ tĩnh}\\ làm cầu nối};
  \node[below=3mm of salt,text=vnCite,font=\tiny,align=center]
       {mượn \emph{mô hình}\\ \emph{đã tiền huấn luyện}};
\end{tikzpicture}}

\vspace{1.0em}
\begin{columns}[T]
\begin{column}{0.94\textwidth}
  \small
  Khởi tạo ngẫu nhiên buộc phần thân học lại từ tín hiệu vô nghĩa
  \cit{deVries2021}. Các phương pháp theo ngữ nghĩa dùng véc-tơ tĩnh để bắc
  cầu; \textbf{SALT} đi xa hơn -- lấy thẳng không gian véc-tơ của một mô hình
  \hl{đã tiền huấn luyện} ở ngôn ngữ đích.
\end{column}
\end{columns}
\end{frame}

% ---------------------------------------------------------------------
\begin{frame}{SALT: tái sử dụng không gian véc-tơ của mô hình ngôn ngữ đích}
\vfill
\centering
\includegraphics[width=\linewidth,height=5.9cm,keepaspectratio]{fig-salt-paper}

\vspace{0.4em}
{\footnotesize\color{vnCite} Nguồn: \citet{Lee2025}}
\vfill
\end{frame}
